\section{Core Mechanics}

% Each subsection below is a slot for one major system. Duplicate
% the pattern --- "what / why / how" --- for every system that
% matters. Keep each system to roughly half a page; if it grows
% larger, split it into its own subsection tree.

\subsection{Player Character}

\paragraph{What it is.}
% One sentence describing the player avatar and its in-engine type.
\textit{[The player is a \ldots implemented as \ldots]}

\paragraph{Why it exists this way.}
% The design intent. What feeling is this character meant to
% produce? What constraints shaped it?
\textit{[\ldots]}

\paragraph{How it is built.}
% The implementation sketch: components, modules, state machine,
% notable interactions with the world.
\textit{[\ldots]}

\subsection{Verbs and Abilities}

% The list of things the player can \emph{do}. For each verb,
% give the input, the in-world effect, and the cost or constraint
% (cooldown, resource, prerequisite). Verbs without a cost tend
% to flatten design space --- be explicit about why if one is free.

\begin{itemize}
    \item \textbf{\textit{[Verb]}} --- \textit{[input / effect / cost].}
    \item \textbf{\textit{[Verb]}} --- \textit{[\ldots].}
\end{itemize}

\subsection{World and Interactables}

% What populates the world the player moves through? Hazards,
% pickups, NPCs, environmental triggers. For each category,
% describe the contract: how it is detected, what it does, what
% it expects from the player.

\textit{[Describe the major categories of interactable here.]}

\subsection{Antagonist / Pressure}

% Whatever pushes the player forward: an enemy, a timer, a
% deteriorating resource, an AI director. Even a puzzle game
% has \emph{some} source of pressure --- name it.

\paragraph{Source of pressure.}
\textit{[\ldots]}

\paragraph{Escalation.}
% How does pressure ramp? Stages, thresholds, triggers?
\textit{[\ldots]}

\paragraph{Resolution.}
% What ends the pressure --- victory, defeat, or a transition?
\textit{[\ldots]}

\subsection{Game States and Flow}

% The high-level state graph: splash, menu, world, pause,
% game-over, credits. Note what owns each state and how
% transitions happen.

\textit{[Describe the top-level state machine.]}

\subsection{Cross-cutting Systems}

% Services that everything depends on: save, audio, input,
% cinematics, post-processing, particles, localization.
% One line each is enough at this stage --- expand later.

\begin{itemize}
    \item \textbf{\textit{[System]}} --- \textit{[one-line role].}
    \item \textbf{\textit{[System]}} --- \textit{[\ldots].}
\end{itemize}

\subsection{Open Design Questions}

% Things that are deliberately unresolved. Naming them here
% prevents them from being forgotten and helps future-you
% recognize when a decision has implicitly been made.

\begin{itemize}
    \item \textit{[Question 1.]}
    \item \textit{[Question 2.]}
\end{itemize}

\newpage
